\chapter{Periodic Fever Syndromes}
\index{Periodic Fever Syndromes}
\label{sec:Periodic Fever Syndromes}

\section{Overview}
Periodic fever syndromes (also called autoinflammatory diseases) are a group of inherited disorders of the innate immune system characterized by recurrent, self-limited episodes of fever and systemic inflammation in the absence of autoantibodies or antigen-specific T cells. They are distinct from autoimmune diseases and arise from dysregulation of the inflammasome and IL-1$\beta$ production (the ``IL-1opathies'').

\subsubsection{Familial Mediterranean Fever}
Familial Mediterranean fever is the most common hereditary periodic fever syndrome, caused by mutations in the \textit{MEFV} gene encoding pyrin, a component of the pyrin inflammasome, most prevalent in Mediterranean populations (Armenian, Turkish, Arab, Jewish---carrier frequency 1:5 in Armenians). This genetic disorder results from specific gene mutations or chromosomal abnormalities. It may be present at birth or manifest later in life depending on the underlying mechanism. Genetic disorders result from abnormalities in an individual's DNA, ranging from single-gene mutations to chromosomal abnormalities. These conditions may be inherited from parents or arise from new mutations. Pattern of inheritance depends on whether the mutation is dominant, recessive, or X-linked. Genetic counseling helps families understand risks and make informed decisions. This condition affects a significant number of people worldwide and can range from mild to severe in presentation. Early recognition and appropriate management are essential for optimal outcomes. A thorough understanding of the underlying mechanisms guides treatment decisions. Patient education and self-management play important roles in long-term care.
\section{Causes}
\section{Risk Factors}

\begin{itemize}[noitemsep]
  \item Genetic predisposition and family history
  \item Advancing age
  \item Lifestyle factors (diet, exercise, smoking, alcohol)
  \item Environmental exposures
  \item Underlying medical conditions
  \item Medication use
\end{itemize}
\section{Signs and Symptoms}

\subsection{Common symptoms}
\begin{itemize}[noitemsep]
  \item Clinical features begin in childhood (90\% before age 20) and consist of recurrent, self-limited attacks (12--72 hours) of high fever, serositis (peritonitis---severe abdominal pain resembling acute abdomen, pleuritis, pericarditis), arthritis (usually monoarticular of the lower extremity), and an erysipelas-like redness of the skin on the lower leg, with the diagnostic hallmark being a dramatic response to colchicine.

  \item Pain or discomfort in the affected area
  \end{itemize}

\subsection{Emergency warning signs}
\begin{itemize}[noitemsep]
  \item Severe or worsening symptoms of periodic fever syndromes require immediate medical evaluation
\end{itemize}
\section{Progression and Stages}
The outlook varies from person to person, with amyloidosis risk reduced with effective anti-IL-1 therapy.

\subsubsection{Cryopyrin-Associated Periodic Syndromes}
Cryopyrin-associated periodic syndromes are a spectrum of autoinflammatory diseases caused by gain-of-function mutations in the \textit{NLRP3} gene encoding cryopyrin (NALP3), a key component of the NLRP3 inflammasome, encompassing three overlapping phenotypes in order of increasing severity: familial cold autoinflammatory syndrome (FCAS---mildest), Muckle-Wells syndrome (MWS---intermediate), and neonatal-onset multisystem inflammatory disease/chronic infantile nervous system cutaneous articular syndrome (NOMID/CINCA---most severe), with an incidence of approximately 1 per 1,000,000. Treatment typically involves with IL-1 inhibitors: anakinra, canakinumab (long-acting, FDA-approved for CAPS), and rilonacept (IL-1 trap), which are dramatically effective, normalizing symptoms and inflammatory markers within hours and preventing AA amyloidosis and hearing loss progression.

\subsubsection{Hyper-IgD Syndrome} The condition may progress through different stages of severity over time. Early stages often involve mild or subtle symptoms that may be overlooked. Moderate stages involve more noticeable symptoms affecting daily function. Advanced stages may involve significant impairment and complications.

\begin{itemize}[noitemsep]
  \item Hyper-IgD syndrome (also known as mevalonate kinase deficiency, MKD) is an autosomal recessive periodic fever syndrome caused by mutations in the \textit{MVK} gene encoding mevalonate kinase, an enzyme in the cholesterol biosynthesis pathway, leading to accumulation of mevalonic acid and activation of the inflammasome, with approximately 300 reported cases, predominantly in Northern European populations (especially Dutch). The outlook is good
  \item Attacks decrease in frequency and severity with age, and life expectancy is normal.
\end{itemize}
\section{Complications}
The most serious long-term complication is amyloidosis (AA type) if untreated.

\begin{itemize}[noitemsep]
  \item Disease progression leading to worsening symptoms
  \item Secondary infections or injuries
  \item Side effects of treatment
  \item Reduced quality of life
  \item Emotional and psychological impact
\end{itemize}
\section{Diagnosis}
Doctors usually diagnose this based on your symptoms and a physical exam (Tel Hashomer criteria) or by genetic confirmation of biallelic \textit{MEFV} mutations.

\begin{itemize}[noitemsep]
  \item Comprehensive medical history and physical examination
  \item Laboratory tests to assess organ function
  \item Imaging studies to visualize affected structures
  \item Specialized testing depending on the affected system
  \item Consultation with relevant specialists
\end{itemize}
\section{Prevention}
\begin{itemize}[noitemsep]
  \item Maintain a healthy lifestyle with balanced diet and exercise
  \item Avoid known risk factors
  \item Attend regular medical check-ups for early detection
  \item Follow recommended screening guidelines
\end{itemize}
\section{Treatment Options}
Treatment typically involves daily colchicine (0.5--2 mg/day), which prevents attacks and amyloidosis in 95\% of patients, with IL-1 inhibitors (anakinra, canakinumab) used for colchicine-resistant FMF. Symptomatic management and supportive care Enzyme replacement therapy for certain conditions Gene therapy (emerging for select conditions) Dietary modifications (e.g., phenylketonuria) Regular monitoring for associated complications Physical and occupational therapy Genetic counseling for family planning Participation in clinical trials for new therapies

\begin{itemize}[noitemsep]
  \item Medications to manage symptoms and address underlying mechanisms
  \item Lifestyle modifications and self-care measures
  \item Physical therapy and rehabilitation
  \item Surgical intervention when indicated
  \item Regular monitoring and follow-up care
\end{itemize}
\section{Living with It}
\begin{itemize}[noitemsep]
  \item Follow your treatment plan as prescribed
  \item Maintain regular follow-up with your healthcare provider
  \item Make necessary lifestyle adjustments
  \item Seek support from family, friends, or support groups
  \item Stay informed about your condition
\end{itemize}
\section{Outlook}
\begin{itemize}[noitemsep]
  \item Most people recover well with colchicine adherence
  \item AA amyloidosis risk is nearly eliminated.
\end{itemize}
